\documentclass[nomath]{harryopo-paper}

\title{基于深度学习的图像识别研究}
\author{张三}
\date{2026年08月05日}


\begin{document}

\maketitle

\section*{一、引言}
\addcontentsline{toc}{section}{一、引言}

\subsection*{1.1 研究背景}
\addcontentsline{toc}{subsection}{1.1 研究背景}

\subsubsection*{1.1.1 深度学习发展}
\addcontentsline{toc}{subsubsection}{1.1.1 深度学习发展}

这是普通文字，\fzht{这是加粗文字}，这是\textit{斜体文字}。

其中损失函数为 L = -Σ y_i log(p_i)，其中 y_i 是真实标签。

\section*{二、实验结果}
\addcontentsline{toc}{section}{二、实验结果}

实验结果如下表所示：

| 方法 | 准确率(%) | 召回率(%) | F1分数 |
| --- | --- | --- | --- |
| ResNet-50 | 94.2 | 92.1 | 93.1 |
| VGG-16 | 91.5 | 89.3 | 90.4 |
| 我们的方法 | 96.8 | 95.2 | 96.0 |

消融实验（含水平合并表头）：

<table><tr><td colspan="2"><p>主干网络</p></td><td><p>性能</p></td></tr><tr><td><p>ResNet</p></td><td><p>ImageNet</p></td><td><p>94.2%</p></td></tr><tr><td><p>VGG</p></td><td><p>ImageNet</p></td><td><p>91.5%</p></td></tr><tr><td><p>EfficientNet</p></td><td><p>ImageNet</p></td><td><p>95.1%</p></td></tr></table>

多数据集对比（含垂直合并）：

<table><tr><td><p>模型</p></td><td><p>数据集</p></td><td><p>准确率</p></td></tr><tr><td rowspan="2"><p>Ours</p></td><td><p>CIFAR-10</p></td><td><p>98.5%</p></td></tr><tr><td><p>CIFAR-100</p></td><td><p>85.3%</p></td></tr><tr><td rowspan="2"><p>Baseline</p></td><td><p>CIFAR-10</p></td><td><p>95.2%</p></td></tr><tr><td><p>CIFAR-100</p></td><td><p>78.9%</p></td></tr></table>

\section*{三、结论}
\addcontentsline{toc}{section}{三、结论}

本文提出的方法在多个数据集上取得了优异的性能。

\end{document}